\section{Experimental Protocol}
\label{sec:protocol}
\subsection{Current validation metrics}
For query $q$ and instance $j$, the existing helper multiplies the sigmoid query logit by the sigmoid concept-presence logit:
\begin{equation}
 s_{qj}=\sigma(z_{qj})\sigma(p_q).
\end{equation}
It accepts instances with $s_{qj}>0.5$, thresholds mask logits at zero, and takes their union:
\begin{equation}
 \widehat G_q=\bigcup_{j:s_{qj}>0.5}\{x:\ell_{qj}(x)>0\}.
\end{equation}
Ground-truth instance masks are area-resampled to the output-mask size and binarized at 0.5 before union. The resulting IoU is averaged across positive queries. It is a polygon-derived, output-resolution score, not native-resolution PNG-mask IoU.

The logged \code{val_recall} is the fraction of positive queries for which at least one instance survives the score threshold. It does not require spatial overlap and is therefore a positive-query response rate rather than lesion detection recall. Similarly, \code{val_false_positive} is the fraction of sampled negative queries with any accepted instance. It is a query-level false-alarm rate, not a pixel false-positive rate. These distinctions must remain explicit when importing training logs into the results section.

\subsection{Final segmentation metrics to add}
The primary final evaluation should use the union of predicted disease regions against the binary PNG reference $G=\{x:Y(x)>0\}$ in the original aligned image frame. With pixel counts TP, FP, and FN, define
\begin{align}
\mathrm{IoU}&=\frac{\mathrm{TP}}{\mathrm{TP}+\mathrm{FP}+\mathrm{FN}},\\
\mathrm{Dice}&=\frac{2\mathrm{TP}}{2\mathrm{TP}+\mathrm{FP}+\mathrm{FN}}.
\end{align}
Pixel precision and recall are $\mathrm{TP}/(\mathrm{TP}+\mathrm{FP})$ and $\mathrm{TP}/(\mathrm{TP}+\mathrm{FN})$. For a nonempty reference and empty prediction, IoU, Dice, and recall are zero; prediction precision is undefined and should be reported as such. Empty references and absent disease classes require an explicit policy rather than an implicit smoothing convention.

Report image-mean IoU and Dice, pooled foreground IoU, and a disease-macro score. For class $c$ with $n_c$ eligible test images,
\begin{equation}
 \mathrm{mIoU}_{\mathrm{disease}}=
 \frac{1}{|\mathcal C_*|}\sum_{c\in\mathcal C_*}
 \frac{1}{n_c}\sum_{i:y_i=c}\mathrm{IoU}_i,
\end{equation}
where $\mathcal C_*$ contains classes with test support. This is an average of image scores within disease, then across diseases. It differs from averaging classwise pooled confusion matrices and from averaging disease and background semantic IoUs.

Dice, precision, disease-macro aggregation, uncertainty estimates, and native-resolution PNG evaluation are not currently produced by the helper. Instance AP may be added only after polygon quality and instance definitions have been resolved. A mask connected component is an operational pseudo-instance and should not be represented as an expert-verified lesion identity.

\subsection{Data preparation and split policy}
Two protocols should be distinguished. The \emph{existing split} protocol preserves the local partition and reports its known duplicate risk. The \emph{quality-controlled} protocol verifies orientation, resolves flagged polygon mismatches, and removes training images belonging to validation or test duplicate clusters. If a cluster spans validation and test, it must also be assigned to one evaluation partition or explicitly excluded. A saved manifest must define these choices, final class support, and exclusions.

The currently available duplicate screening can initialize this process, but its threshold-based clusters need review before a final manifest is frozen. No corrected manifest has been generated by this report. Dataset tables describe the available files and current builder, not an already cleaned benchmark.

Hyperparameters, stopping criteria, and thresholds should be selected on validation only. The test partition should be evaluated after those decisions are fixed. Include all 1,561 v7 test PNG pairs for semantic evaluation unless the chosen quality policy explicitly removes samples. Document every exclusion, including any difference from the builder's 1,560 polygon-supported images.

\subsection{Comparisons and ablations}
The primary comparison is pretrained SAM 3 without adaptation versus the default frozen-encoder adaptation, using identical disease prompts, preprocessing, and scoring. A full-adaptation configuration can test whether updating representations adds value when resources permit. Further controlled ablations can vary positive-only versus negative-query training, same-plant versus other-plant negatives, and original versus quality-controlled supervision.

U-Net and DeepLabv3+ are optional future semantic baselines. Their output task and use of disease labels must be specified to make a comparison interpretable. A generic lesion model that receives no disease name has a different input contract from a model given the correct disease phrase. None of these prospective baselines or ablations is claimed as completed.

Use multiple training seeds if feasible and report their mean and variation. Resampling uncertainty should respect duplicate clusters rather than treating near-identical images as independent observations. Also analyze small-coverage images and rare diseases separately. Prompted localization is not a substitute for a disease-classification experiment, so conclusions should remain conditional on the supplied prompt.

\subsection{Reproducibility records}
A final run should preserve the pretrained checkpoint identifier and hash, source revision, resolved arguments, dataset manifest, prompt vocabulary, seed, software versions, GPU model, optimizer updates, and checkpoint-selection rule. Record peak allocated GPU memory and latency with input size, precision, batch size, and synchronization policy. All these runtime measurements remain pending.

The report evidence snapshot records source and dataset hashes and the vendored SAM 3 revision. The existing script produces JSONL training logs and adaptation checkpoints. Resuming additionally requires a compatible base model and optimizer state. The main entry point currently runs training and validation; a standalone test evaluator must be supplied for the final protocol above.
